\section{Modular stability and completeness}\label{sec:modular-stability}

Assume now that the plane is \(\PG(2,q)\), where \(q=p^e\).  Dualizing the
set of \(n\)-secants gives a point set \(\mathcal L^*\) in the dual plane.
The intersection of \(\mathcal L^*\) with the line dual to a point
\(X\) has size \(\rho_n(X)\) or \(\sigma_n(X)\), according as
\(X\in\mathcal K\) or \(X\notin\mathcal K\).  At a rational equality
parameter from Section~\ref{sec:rational-parameters}, the two limiting
intersection numbers are
\[
 \lambda=uv,\qquad a=(u+1)(v+1),\qquad a-\lambda=u+v+1.
\]
When \(u+v+1\) is a power of \(p\), these two numbers are congruent modulo
\(p\).

Sz\H{o}nyi and Weiner prove that a multiset in \(\PG(2,p^e)\) for which only a
small number of lines fail to have one fixed intersection residue can be
changed at few points into a multiset having that residue on every line
\cite[Theorem~1.2]{SWmodp}.  We use their theorem in the range where the
number of exceptional lines is \(O(q)\), so the number of changed support
points is bounded while \(q\) tends to infinity.

The following observation is where completeness re-enters the modular
argument.

\begin{lemma}[Contribution of the repair support]
\label{lem:repair-support}
\coverage{absent}
Fix an ordered factorization \(\lambda=uv\) and suppose
\(d=u+v+1\) is a power of the characteristic.  Let
\[
 q=dm,\qquad n=(u+1)m+1,\qquad
 k=udm^2+cm+o(m),\qquad
 T=ud(v+1)m+h+o(1),
\]
where \(h\in\mathbb Z\).  Suppose every external point is on at least
\(\lambda\) \(n\)-secants, and suppose \(\mathcal L^*\) can be changed at
\(r\) distinct support points into a multiset meeting every line in
\(\lambda\pmod p\) points.  If \(r\) is fixed and nonzero, then
\begin{equation}\label{eq:repair-support-bound}
 c\ge L(h)+\frac ru.
\end{equation}
\end{lemma}

\begin{proof}
Dualize the \(r\) support points to \(r\) lines of the original plane.
Away from their pairwise intersections, a point on one of these lines has
\(n\)-secant multiplicity different from \(\lambda\pmod p\).  If it is
external, its multiplicity is at least \(\lambda+1\).  A support line contains
at most \(n\) points of \(\mathcal K\), and the union loses only
\(O(r^2)\) points at pairwise intersections.  Hence the total external excess
over multiplicity \(\lambda\) is at least
\begin{equation}\label{eq:repair-excess-lower}
 r(q+1-n)-O(r^2)=rvm-O(r^2).
\end{equation}
On the other hand, the incidence sum gives the exact first-order expression
\begin{equation}\label{eq:repair-excess-exact}
 (q+1-n)T-\lambda(q^2+q+1-k)
 =uv\bigl(c-L(h)\bigr)m+o(m).
\end{equation}
Comparing \eqref{eq:repair-excess-lower} and
\eqref{eq:repair-excess-exact} proves \eqref{eq:repair-support-bound}.
\end{proof}

There is a useful restriction when \(r=1\).  Summing the line intersections
of an exact \(\lambda\pmod p\) multiset over a pencil shows that its total
multiplicity is \(\lambda\pmod p\).  Along a characteristic-compatible
family the leading term of \(T\) is divisible by \(p\), so the correction has
total multiplicity \(h-\lambda\pmod p\).  Consequently, if
\begin{equation}\label{eq:support-one-obstruction}
 h\equiv\lambda\pmod p,
\end{equation}
then a nonzero correction cannot have support one.

We apply this first to ordinary completeness.  Write \(q=3m\),
\(n=2m+1\), and suppose
\begin{equation}\label{eq:characteristic-three-asymptotic}
 k=3m^2+cm+o(m),\qquad T=6m+h+o(1).
\end{equation}
The coverage and pair-count inequalities from Sections~\ref{sec:n-secants}
and \ref{sec:rational-parameters} become
\begin{equation}\label{eq:characteristic-three-linear-bounds}
 c+h\ge3,\qquad 5c-h\ge18.
\end{equation}

\begin{lemma}[The uncorrected case is impossible]
\label{lem:no-uncorrected-characteristic-three}
\coverage{absent}
Under \eqref{eq:characteristic-three-asymptotic}, if
\(c<4\) and every line meets \(\mathcal L^*\) in \(1\pmod3\) points, then
the exact pair identity \eqref{eq:n-secant-pairs} cannot hold.
\end{lemma}

\begin{proof}
Near this equality parameter the balanced internal multiplicities are
\(3,4\), and the balanced external multiplicities are \(1,2\).  Let \(D\)
be the total internal deficit below four and \(E\) the total external excess
above one.  The incidence sums give
\begin{align}
 D&=(4c-2h-6)m+o(m),\label{eq:D-characteristic-three}\\
 E&=(c+h-3)m+o(m).\label{eq:E-characteristic-three}
\end{align}
Without the congruence restriction, the available pair-count slack is
\begin{equation}\label{eq:pair-slack-characteristic-three}
 (5c-h-18)m+o(m).
\end{equation}
If every multiplicity is \(1\pmod3\), the internal values nearest four are
\(1,4\), so the minimum pair count rises by \(D\); the external values
nearest one are \(1,4\), so it rises by \(E\).  Feasibility would therefore
require
\[
 5c-h-18\ge \frac{D+E}{m}+o(1)=5c-h-9+o(1),
\]
which is impossible.
\end{proof}

\begin{proof}[Proof of Theorem~\ref{thm:characteristic-three}]
Suppose the assertion fails.  For some \(\varepsilon>0\), pass to a sequence
with
\[
 k\le3m^2+(4-3\varepsilon)m.
\]
The leading real inequality has a unique equality point at \(T/m=6\), so
\(T=6m+o(m)\).  The first-order form of the integer pair condition then gives
\(T-6m=O(1)\); pass to a subsequence on which this difference is a fixed
integer \(h\).

The total deviation from the two balanced degree distributions is \(O(m)\).
Hence only \(O(q)\) dual lines fail to meet \(\mathcal L^*\) in
\(1\pmod3\) points.  Sz\H{o}nyi--Weiner's repair theorem changes
\(\mathcal L^*\) at a fixed number \(r\) of support points to obtain an exact
\(1\pmod3\) multiset.  Lemma~\ref{lem:no-uncorrected-characteristic-three}
excludes \(r=0\).  Lemma~\ref{lem:repair-support} gives
\[
 c\ge L(h)+r=3-h+r.
\]
By \eqref{eq:support-one-obstruction}, \(r\ge2\) when
\(h\equiv1\pmod3\), and \(r\ge1\) otherwise.  Together with the pair-count
bound \(c\ge(18+h)/5\), this yields
\begin{equation}\label{eq:characteristic-three-minimum}
 c\ge
 \min_{h\in\mathbb Z}
 \max\left\{\frac{18+h}{5},,3-h+r_{\min}(h)\right\}=4,
\end{equation}
where \(r_{\min}(h)=2\) for \(h\equiv1\pmod3\) and is one otherwise.  This
contradicts the chosen sequence and proves
\eqref{eq:characteristic-three-bound}.
\imports{SWmodp:repair}
\uses{lem:repair-support,lem:no-uncorrected-characteristic-three,cor:n-secant-feasibility}
\end{proof}
